% Candidate tax events outside the SSB category.
% Raw material from a research agent pass on state tax variation
% 2010--2023. Verified against NCSL State Tax Actions, CMS CCIIO,
% Justia state codes, and state DOR sources cited inline.
% Not a complete 50-state coded panel --- a verified candidate list.

\section*{Non-SSB tax-change candidates, 2010--2023}

\subsection*{Headline: where the variation lives}

\begin{itemize}
  \item \textbf{Almost nothing in general state insurance premium
        taxes.} Zero clean $\geq$50bp P/C or life/health rate
        changes found in 2010--2023. Near-misses are 25--30bp.
  \item \textbf{Indiana financial institutions tax} is the
        cleanest banking variation. It's a tax \emph{cut}
        rather than an imposition, but if the underlying
        mechanism is symmetric (margin shifts $\to$ ads
        shift), cuts and hikes are the same experiment with
        opposite-sign predicted ad responses.
  \item \textbf{State health-insurer assessments} are the
        richest category. Most were designed to backfill the
        federal ACA 9010 fee, so they bundle with reinsurance
        and federal-waiver policy.
  \item \textbf{Hospital provider taxes} are real but bundled
        with Medicaid financing and supplemental-payment
        mechanics. Not pure tax shocks.
\end{itemize}

\subsection*{Insurance premium tax rate changes $\geq$50bp}

None found. Near-misses for reference:
\begin{itemize}
  \item Connecticut, Jan 2018: 1.75\% $\to$ 1.5\% (−25bp).
  \item Arizona, Jan 2016 phase-in: 2.0\% $\to$ 1.7\% over 11
        years (−30bp total).
  \item Alaska, Jan 2016: life premiums over \$100k bracket
        0.10\% $\to$ 0.08\% (−2bp).
\end{itemize}
Source: NCSL State Tax Actions.

\subsection*{Bank franchise / financial institution tax}

\textbf{Indiana FIT phase-down.} Best clean banking
candidate.
\begin{itemize}
  \item CY2013 8.5\% $\to$ CY2014 8.0\% $\to$ CY2015 7.5\%
        $\to$ CY2016 7.0\% $\to$ CY2017 6.5\% $\to$ CY2019
        6.25\% $\to$ CY2020 6.0\% $\to$ CY2021 5.5\% $\to$
        CY2022 5.0\% $\to$ after CY2022 4.9\%.
  \item Qualifying $\geq$50bp steps: 2014, 2015, 2016, 2017,
        2021, 2022. Total drop $-$360bp ($\sim$$-$42\%).
  \item Caveat: it's a franchise/income tax, not a per-unit
        excise. No consumer-side price pass-through
        mechanism, so the pass-through-heterogeneity slice
        doesn't apply. Margin-channel logic still works
        (lower tax $\to$ higher after-tax margin $\to$ more
        ads), with predicted sign flipped vs.\ SSB case.
\end{itemize}

\textbf{North Carolina bank privilege tax repeal} (2016, via
S.L. 2015-241/268). Base removal, not a clean rate change.
Use as a repeal event if needed.

\textbf{Pennsylvania Bank Shares Tax}, Jan 2017: 0.89\%
$\to$ 0.95\% (+6bp). Subthreshold.

\subsection*{State health-insurer assessments (richest
  category)}

These are premium-based, so there \emph{is} a consumer-side
price mechanism. Almost all bundle with reinsurance / 1332
waiver / ACA 9010 replacement design, which contaminates
identification.

\begin{itemize}
  \item \textbf{New Mexico}, Jan 1, 2022: health insurance
        premium surtax 1\% $\to$ 3.75\% (+275bp). Funds the
        Health Care Affordability Fund. Statutory; clean
        effective date.
  \item \textbf{New Jersey}, 2021 first collection: 0 $\to$
        2.5\% (+250bp). Health Insurance Affordability Fund
        assessment on net written premiums of certain health
        and dental benefit issuers.
  \item \textbf{Maryland}, 2019: 0 $\to$ 2.75\% (+275bp);
        2020--2023: 1\% (still +100bp vs.\ no fee). Tied to
        Maryland's State Reinsurance Program and 1332
        waiver --- bundled.
  \item \textbf{Delaware}, 2020: 0 $\to$ 2.75\% in years
        federal HIPF/9010 is waived; 1\% in years HIPF is
        assessed. Bundled with individual-market reinsurance.
  \item \textbf{Colorado}, 2021: 0 $\to$ 1.15\% (nonprofit
        carriers); 0 $\to$ 2.1\% (for-profit). Bundled with
        reinsurance, affordability subsidies, and a separate
        hospital special assessment (\$20m/year 2022--2023)
        which the statute explicitly bars from being passed
        through to consumers.
  \item \textbf{Oregon}, 2018 / Jan 1, 2020: 0 $\to$ 1.5\%
        (HB 2391, 2017), then 1.5\% $\to$ 2.0\% (HB 2010,
        2019). 2020 step is exactly +50bp but also expanded
        base to stop-loss insurance --- rate and base move
        together.
  \item \textbf{Montana}, 2020: 0 $\to$ 1.2\% on major
        medical premiums (+120bp). Lower confidence ---
        sourced from CMS CCIIO but Montana statute not
        located in the research pass.
\end{itemize}

\subsection*{Hospital / provider assessments}

Real tax shocks but heavily bundled with Medicaid financing
and supplemental-payment mechanics. Use only if the
identifying story is about provider-side ads net of
reimbursement changes.

\begin{itemize}
  \item \textbf{Tennessee}, July 2010 initial 3.52\%; raised
        to 4.52\% in subsequent annual renewals (+100bp).
        Statute restricts pass-through: covered hospitals
        may not increase charges or add a surcharge because
        of the assessment.
  \item \textbf{Wyoming}, July 1, 2016: new private hospital
        assessment, up to 2\% of net patient revenue in year
        1. Bundled with Upper Payment Limit / Medicaid
        matching.
  \item \textbf{Connecticut}, effective Nov 21, 2017 (quarters
        from July 1, 2017): hospital tax formula/base
        change. Inpatient 6\% of FY2016 audited inpatient
        net revenue; outpatient 12.3325\% declining. \$343.9m
        revenue effect (NCSL). Code as base/formula change,
        not a clean rate increase.
  \item \textbf{Louisiana}, Aug 1, 2018: emergency ground
        transportation provider fee 3.5\% $\to$ 6\%
        (+250bp). Not a hospital bed tax; include only if
        category 3 is broadened.
\end{itemize}

Subthreshold: Rhode Island hospital license fee 5.65\% $\to$
5.86\% FY2018 (+21bp); West Virginia acute-care hospital tax
0.74\% $\to$ 0.75\% (+1bp).

\subsection*{Ranked candidates for the advertising question}

Best fits for a soda-tax-style ad response test, in order:

\begin{enumerate}
  \item \textbf{New Mexico 2022 surtax (1\% $\to$ 3.75\%).}
        Large, premium-based, clean statutory effective date.
  \item \textbf{New Jersey 2021 HIA (0 $\to$ 2.5\%).} Large
        new imposition; good for health-insurer advertising.
  \item \textbf{Maryland / Delaware 2019--2020 assessments.}
        Large shocks but heavily bundled with reinsurance.
  \item \textbf{Oregon 2020 step (1.5\% $\to$ 2.0\%).}
        Exactly 50bp; base also expanded. Mixed event.
  \item \textbf{Indiana FIT phase-down.} Cleanest statistically;
        opposite-sign predicted response, no consumer price
        mechanism (so no pass-through-heterogeneity slice).
\end{enumerate}

For hospitals (Tennessee, Wyoming): real tax shocks but the
Medicaid-financing bundle is severe.

\subsection*{Source notes}

NCSL State Tax Actions database; CMS CCIIO 1332 data brief
(June 2020); Justia state-code pages for IN, NM, CO, CT,
WY; Tennessee TennCare and Capitol bill text; Oregon DFR;
NJDOBI bulletin. URLs in the original research-agent
output (saved separately, ask if needed).
